\NSPage{121}%
\begin{EESegment}{EE-TGT-P121-001}
\subsection{Energy and comparison.}\label{sec:10.3}
This subsection starts by proving that the smooth compactly supported force
\NSi{NS-F-P121-001}{f} from Lemma~\ref{lem:10.3} gives an energy bound for the
constructed velocity \NSi{NS-F-P121-002}{u} that stays uniform up to time one
(Lemma~\ref{lem:10.4}). Next it proves that any smooth solution with the same force,
zero initial data, and bounded kinetic energy must agree with that velocity on every
shorter time interval (Lemma~\ref{lem:10.5}). This comparison lets the proof carry the
local velocity growth in \eqref{eq:3.6} over to any supposed global smooth solution. In
these energy estimates on all of space, \NSi{NS-F-P121-003}{\|\cdot\|_p} is the spatial
\NSi{NS-F-P121-004}{L^p(\R^3)} norm, and every spatial integral with no domain written
beside it is over \NSi{NS-F-P121-005}{\R^3}.
\end{EESegment}

\begin{lemma}\label{lem:10.4}
\begin{EESegment}{EE-TGT-P121-002}
Set \NSi{NS-F-P121-006}{\mathcal F(t)=\int_0^t\|f(s)\|_2\,ds} for
\NSi{NS-F-P121-007}{0\leq t\leq 1}. Then
\NSi{NS-F-P121-008}{\mathcal F(1)<\infty}, and
\NSFormula{NS-F-P121-009}{display}{10.13}
\begin{equation}
  \|u(t)\|_2^2+2\int_0^t\|\nabla u(s)\|_2^2\,ds
  \leq \mathcal F(t)^2
  \qquad (0\leq t<1).
\tag{10.13}\label{eq:10.13}
\end{equation}
In particular, the kinetic energy stays uniformly bounded, and the total dissipation on
\NSi{NS-F-P121-010}{[0,1)} is finite.
\end{EESegment}
\end{lemma}

\begin{proof}
\begin{EESegment}{EE-TGT-P121-003}
On any interval \NSi{NS-F-P121-011}{[0,T]} with \NSi{NS-F-P121-012}{T<1}, the
localized velocity and pressure are smooth and have support in one fixed compact spatial
set. Taking the inner product of the equation with \NSi{NS-F-P121-013}{u} and
integrating in space gives
\NSFormula{NS-F-P121-014}{display}{10.14}
\begin{equation}
  \frac12\frac{d}{dt}\|u(t)\|_2^2+\|\nabla u(t)\|_2^2
  =\langle f(t),u(t)\rangle.
\tag{10.14}\label{eq:10.14}
\end{equation}
For \NSi{NS-F-P121-015}{\delta>0}, divide this identity by
\NSi{NS-F-P121-016}{(\|u(t)\|_2^2+\delta^2)^{1/2}}, drop the nonnegative dissipation
term, and apply Cauchy--Schwarz. Integrate from the zero initial data and then let
\NSi{NS-F-P121-017}{\delta\downarrow0}. The result is
\NSi{NS-F-P121-018}{\|u(t)\|_2\leq\mathcal F(t)}. Put this bound back into
\eqref{eq:10.14} and integrate once more to get
\NSFormula{NS-F-P121-019}{display}{none}
\[
  \|u(t)\|_2^2+2\int_0^t\|\nabla u(s)\|_2^2\,ds
  \leq 2\int_0^t\mathcal F'(s)\mathcal F(s)\,ds
  =\mathcal F(t)^2.
\]
Because the force is smooth and compactly supported,
\NSi{NS-F-P121-020}{\mathcal F(1)<\infty}. Monotone convergence then shows that the
integrated dissipation on \NSi{NS-F-P121-021}{[0,1)} is finite, using only the bounds
already proved for \NSi{NS-F-P121-022}{t<1}.
\end{EESegment}
\end{proof}

\begin{EESegment}{EE-TGT-P121-004}
The comparison argument only needs smoothness and a uniform spatial
\NSi{NS-F-P121-023}{L^2} bound on \NSi{NS-F-P121-024}{[0,T]}, where
\NSi{NS-F-P121-025}{T<1}. These assumptions put no restriction on how the spatial
derivatives grow at infinity. The equation will first be used to find the pressure gradient,
and only then will the proof take the limit in the localized energy identity.
\end{EESegment}

\begin{lemma}\label{lem:10.5}
\begin{EESegment}{EE-TGT-P121-005}
Fix \NSi{NS-F-P121-026}{T<1}. Suppose \NSi{NS-F-P121-027}{v,P} is a smooth solution
of \eqref{eq:1.1} with viscosity one on
\NSi{NS-F-P121-028}{\R^3\times[0,T]}. Suppose it has the force
\NSi{NS-F-P121-029}{f} from Lemma~\ref{lem:10.3}, zero initial velocity, and
\NSi{NS-F-P121-030}{v\in L^\infty([0,T];L^2(\R^3))}. Then
\NSi{NS-F-P121-031}{v=u} on that whole interval, where
\NSi{NS-F-P121-032}{u} is the localized velocity from Proposition~\ref{prop:10.1}.
\end{EESegment}
\end{lemma}

\begin{proof}
\begin{EESegment}{EE-TGT-P121-006}
Take the difference of the two solutions and estimate it on balls whose radii grow without
bound. The pressure term needs its own estimate because nothing has been assumed about the
spatial growth of \NSi{NS-F-P121-033}{P}. Each constant written as
\NSi{NS-F-P121-034}{C_T} below may depend on \NSi{NS-F-P121-035}{T} and on the
stated bounds for the fields, but it does not depend on the cutoff radius. Set
\NSFormula{NS-F-P121-036}{display}{none}
\[
  w=v-u,\qquad
  \pi=P-p,\qquad
  g_{ij}=v_iv_j-u_iu_j=w_iw_j+w_iu_j+u_iw_j.
\]
The convention is
\NSi{NS-F-P121-037}{(\operatorname{div}g)_i=\sum_j\partial_jg_{ij}}.
On \NSi{NS-F-P121-038}{[0,T]},
\NSi{NS-F-P121-039}{\|w(t)\|_2\leq C_T} and
\NSi{NS-F-P121-040}{\sum_{i,j}\|g_{ij}(t)\|_1\leq C_T}.
The same force appears in both equations, so it cancels and leaves
\NSFormula{NS-F-P121-041}{display}{10.15}
\begin{equation}
  \partial_t w+(v\cdot\nabla)w+(w\cdot\nabla)u
  =\Delta w-\nabla\pi,\qquad \operatorname{div}w=0.
\tag{10.15}\label{eq:10.15}
\end{equation}
\end{EESegment}

\begin{EESegment}{EE-TGT-P121-007}
\paragraph{\normalfont\itshape Pressure gradient.}
Let \NSi{NS-F-P121-042}{R_i} be the Riesz transform whose Fourier multiplier is
\NSi{NS-F-P121-043}{i\xi_i/|\xi|}, and define the pressure as a distribution, a continuous linear
rule that assigns a number to each smooth compactly supported test function, by
\NSFormula{NS-F-P121-044}{display}{10.16}
\begin{equation}
  \pi_*=\sum_{i,j=1}^3 R_iR_jg_{ij}.
\tag{10.16}\label{eq:10.16}
\end{equation}
\end{EESegment}

\NSPage{122}%
\begin{EESegment}{EE-TGT-P122-001}
The multiplier for this operator is
\NSi{NS-F-P122-001}{-\xi_i\xi_j/|\xi|^2}; its value at zero does not matter. The
Fourier transform of an \NSi{NS-F-P122-002}{L^1} function is bounded, so
\NSi{NS-F-P122-003}{\pi_*} is uniformly in
\NSi{NS-F-P122-004}{H^{-s}(\R^3)} for every \NSi{NS-F-P122-005}{s>3/2}. More
exactly, its squared norm is at most a constant times
\NSi{NS-F-P122-006}{\|g\|_1^2\int_{\R^3}(1+|\xi|^2)^{-s}\,d\xi}. The multiplier also
gives
\NSFormula{NS-F-P122-007}{display}{none}
\[
  \Delta\pi_*=-\sum_{i,j}\partial_i\partial_jg_{ij}.
\]
The next step is to prove that \NSi{NS-F-P122-008}{\nabla\pi=\nabla\pi_*} in the
space--time interior. Take \NSi{NS-F-P122-009}{a\in C_c^\infty(0,T)}. Writing
\eqref{eq:10.15} in conservative form gives
\NSFormula{NS-F-P122-010}{display}{none}
\[
  \int_0^T a\nabla\pi\,dt
  =\Delta\int_0^T aw\,dt+\int_0^T a'w\,dt
   -\operatorname{div}\int_0^T ag\,dt.
\]
The three terms on the right belong, in that order, to
\NSi{NS-F-P122-011}{H^{-2}}, \NSi{NS-F-P122-012}{L^2}, and
\NSi{NS-F-P122-013}{H^{-3}}. The last one has this property because
\NSi{NS-F-P122-014}{g\in L_t^\infty L_x^1}, and the Fourier estimate above can be
used with \NSi{NS-F-P122-015}{s=2}. Also,
\NSi{NS-F-P122-016}{\pi_*\in L_t^\infty H_x^{-2}}, so
\NSFormula{NS-F-P122-017}{display}{none}
\[
  H_a:=\int_0^T a(\nabla\pi-\nabla\pi_*)\,dt\in H^{-3}(\R^3).
\]
Taking the divergence of the difference equation gives
\NSi{NS-F-P122-018}{\Delta\pi=-\sum_{i,j}\partial_i\partial_jg_{ij}=\Delta\pi_*}.
It follows that \NSi{NS-F-P122-019}{\Delta H_a=0}. The Fourier transform of
\NSi{NS-F-P122-020}{H_a} is a weighted \NSi{NS-F-P122-021}{L^2} function supported
at the single point \NSi{NS-F-P122-022}{\{0\}}, so it must be zero. Allowing a spatial
test function as well proves the stated equality. The equation has now determined the
pressure gradient under these assumptions without putting any restriction on the spatial
growth of \NSi{NS-F-P122-023}{P}.
\end{EESegment}

\begin{EESegment}{EE-TGT-P122-002}
\paragraph{\normalfont\itshape Pressure flux.}
Next, bound the pressure contribution on the boundary of an expanding ball by the weighted
dissipation inside that ball. Start with
\NSi{NS-F-P122-024}{\int\pi w\cdot\nabla\chi_R}. It will be estimated by powers of a
weighted \NSi{NS-F-P122-025}{L^6} norm of \NSi{NS-F-P122-026}{w}, with powers low
enough for the dissipation to absorb them. Choose a smooth, compactly supported function
\NSi{NS-F-P122-027}{0\leq\varphi\leq1} that equals one on the unit ball. For
\NSi{NS-F-P122-028}{R\geq1}, write
\NSi{NS-F-P122-029}{\varphi_R(x)=\varphi(x/R)}, set
\NSi{NS-F-P122-030}{\chi_R=\varphi_R^8}, and define
\NSFormula{NS-F-P122-031}{display}{none}
\[
  E_R=\int\chi_R|w|^2,\qquad
  A_R=\left(\int\chi_R|\nabla w|^2\right)^{1/2},\qquad
  B_R=\|\varphi_R^4w\|_6.
\]
The Sobolev inequality and the product rule give
\NSFormula{NS-F-P122-032}{display}{10.17}
\begin{equation}
  B_R\leq C\|\nabla(\varphi_R^4w)\|_2
  \leq C(A_R+R^{-1}\|w\|_2).
\tag{10.17}\label{eq:10.17}
\end{equation}
Use the standard fact that Riesz transforms are bounded on \NSi{NS-F-P122-033}{L^p}
when \NSi{NS-F-P122-034}{1<p<\infty}; see \cite{ref-20}. If
\NSi{NS-F-P122-035}{a} is a scalar function and \NSi{NS-F-P122-036}{\mathcal T} is a
linear operator, the commutator
\NSi{NS-F-P122-037}{[a,\mathcal T]g=a(\mathcal Tg)-\mathcal T(ag)} measures the
difference between multiplying before and after applying the operator. Multiplying
\eqref{eq:10.16} by \NSi{NS-F-P122-038}{\varphi_R^4} gives
\NSFormula{NS-F-P122-039}{display}{10.18}
\begin{equation}
  \varphi_R^4\pi_*
  =\sum_{i,j}R_iR_j(\varphi_R^4g_{ij})
   +\sum_{i,j}[\varphi_R^4,R_iR_j]g_{ij}.
\tag{10.18}\label{eq:10.18}
\end{equation}
The first sum has \NSi{NS-F-P122-040}{L^{3/2}} norm no larger than
\NSi{NS-F-P122-041}{C_T(B_R+1)}, because
\NSFormula{NS-F-P122-042}{display}{none}
\[
  \|\varphi_R^4w_iw_j\|_{3/2}\leq B_R\|w\|_2,\qquad
  \|\varphi_R^4w_iu_j\|_{3/2}\leq\|w\|_2\|u\|_6.
\]
The nonlocal kernel of \NSi{NS-F-P122-043}{R_iR_j} is bounded by
\NSi{NS-F-P122-044}{C|x-y|^{-3}}. Any multiple of the identity in that operator
cancels from the commutator, so the absolute value of the commutator kernel is bounded by
\NSFormula{NS-F-P122-045}{display}{none}
\[
  K_R(x-y)=C|x-y|^{-3}\min\{|x-y|/R,1\}.
\]
Integrating directly in the radial variable gives
\NSFormula{NS-F-P122-046}{display}{none}
\[
  \|K_R\|_{4/3}^{4/3}
  \leq C\left[
    R^{-4/3}\int_0^R r^{-2/3}\,dr+\int_R^\infty r^{-2}\,dr
  \right]
  \leq CR^{-1}.
\]
\end{EESegment}

\NSPage{123}%
\begin{EESegment}{EE-TGT-P123-001}
Young's convolution inequality now bounds the second sum in \eqref{eq:10.18} in
\NSi{NS-F-P123-001}{L^{4/3}} by \NSi{NS-F-P123-002}{C_TR^{-3/4}}. These identities
can first be proved for compact smooth cutoffs of \NSi{NS-F-P123-003}{g}, then passed to
\NSi{NS-F-P123-004}{g} by convergence in \NSi{NS-F-P123-005}{L^1}, the uniform
commutator bound, and convergence of the Riesz transforms in
\NSi{NS-F-P123-006}{H^{-s}}. In particular,
\NSi{NS-F-P123-007}{\pi_*} is locally integrable in both space and time.
\end{EESegment}

\begin{EESegment}{EE-TGT-P123-002}
Test the equality of the pressure gradients against the compact vector field
\NSi{NS-F-P123-008}{a(t)\chi_Rw}. This gives
\NSFormula{NS-F-P123-009}{display}{none}
\[
  \int\pi w\cdot\nabla\chi_R\,dx
  =\int\pi_*w\cdot\nabla\chi_R\,dx
  \qquad\text{for almost every }t\in(0,T).
\]
Because \NSi{NS-F-P123-010}{|\nabla\chi_R|\leq CR^{-1}\varphi_R^7}, the two sums in
\eqref{eq:10.18} are paired with
\NSi{NS-F-P123-011}{CR^{-1}\varphi_R^3|w|}. Interpolation gives
\NSFormula{NS-F-P123-012}{display}{none}
\[
  \|\varphi_R^3w\|_3\leq\|\varphi_R^2w\|_3
  \leq B_R^{1/2}\|w\|_2^{1/2},
  \qquad
  \|\varphi_R^3w\|_4\leq B_R^{3/4}\|w\|_2^{1/4}.
\]
It follows that
\NSFormula{NS-F-P123-013}{display}{10.19}
\begin{equation}
  \left|\int\pi w\cdot\nabla\chi_R\right|
  \leq C_TR^{-1}\left[(B_R+1)B_R^{1/2}+R^{-3/4}B_R^{3/4}\right].
\tag{10.19}\label{eq:10.19}
\end{equation}
\end{EESegment}

\begin{EESegment}{EE-TGT-P123-003}
\paragraph{\normalfont\itshape Difference energy.}
The pressure-flux bound can now be used in the energy estimate for the difference. Every
integral has compact support, so smoothness is enough to justify the calculation. Pairing
\eqref{eq:10.15} with \NSi{NS-F-P123-014}{\chi_Rw} gives
\NSFormula{NS-F-P123-015}{display}{none}
\[
\begin{aligned}
  \frac12E_R'+A_R^2
  ={}&-\int\chi_R(w\cdot\nabla)u\cdot w
      +\frac12\int|w|^2\Delta\chi_R\\
     &+\frac12\int|w|^2v\cdot\nabla\chi_R
      +\int\pi w\cdot\nabla\chi_R.
\end{aligned}
\]
Choose \NSi{NS-F-P123-016}{R} large enough that
\NSi{NS-F-P123-017}{\chi_R=1} on a neighborhood of
\NSi{NS-F-P123-018}{\operatorname{supp}u}. Then
\NSi{NS-F-P123-019}{v=w} on
\NSi{NS-F-P123-020}{\operatorname{supp}\nabla\chi_R}, and the transport flux is no
larger than
\NSFormula{NS-F-P123-021}{display}{none}
\[
  CR^{-1}\int\varphi_R^6|w|^3\leq C_TR^{-1}B_R^{3/2}.
\]
The last bound follows from
\NSi{NS-F-P123-022}{\varphi_R^2|w|=(\varphi_R^4|w|)^{1/2}|w|^{1/2}} and Hölder's
inequality. The Laplacian term is at most \NSi{NS-F-P123-023}{C_TR^{-2}}. Substitute
\eqref{eq:10.17} into \eqref{eq:10.19}. Every power of
\NSi{NS-F-P123-024}{A_R} that appears in these flux bounds is then no larger than
\NSi{NS-F-P123-025}{3/2<2}, so Young's inequality gives
\NSFormula{NS-F-P123-026}{display}{none}
\[
  \frac12E_R'+\frac12A_R^2
  \leq\|\nabla u\|_\infty E_R+\frac{C_T}{R}
  \qquad\text{for almost every }t\in(0,T).
\]
For example,
\NSi{NS-F-P123-027}{R^{-1}A_R^{3/2}\leq\delta A_R^2+C_\delta R^{-4}}; the terms
with powers \NSi{NS-F-P123-028}{1/2} and \NSi{NS-F-P123-029}{3/4} obey the same
required bound. The coefficient \NSi{NS-F-P123-030}{\|\nabla u\|_\infty} is bounded on
\NSi{NS-F-P123-031}{[0,T]}, and \NSi{NS-F-P123-032}{E_R(0)=0}. Gronwall's inequality
now gives
\NSFormula{NS-F-P123-033}{display}{none}
\[
  E_R(t)\leq\frac{2C_T}{R}\int_0^t
  \exp\left(2\int_s^t\|\nabla u(a)\|_\infty\,da\right)\,ds
  \leq\frac{C_T'}{R}.
\]
For every fixed ball, \NSi{NS-F-P123-034}{\chi_R=1} once
\NSi{NS-F-P123-035}{R} is large enough. Letting \NSi{NS-F-P123-036}{R\to\infty}
therefore gives \NSi{NS-F-P123-037}{w=0} throughout
\NSi{NS-F-P123-038}{[0,T]}.
\end{EESegment}
\end{proof}

\begin{EESegment}{EE-TGT-P123-004}
\subsection{Growth, lifespan, and viscosity.}\label{sec:10.4}
The force extension and energy estimate for the localized fields
\NSi{NS-F-P123-039}{(u,p)} are now in place (Lemma~\ref{lem:10.3},
\eqref{eq:10.13}). This subsection uses the inner growth asymptotic \eqref{eq:3.6} from
Theorem~\ref{thm:3.1} to find the maximal classical lifespan. Together with the comparison
lemma, Lemma~\ref{lem:10.5}, the same formula rules out a global smooth solution with
bounded energy. A spatial rescaling then gives the result for every fixed positive viscosity
without changing the time of the singularity.
\end{EESegment}

\NSPage{124}%
\begin{proof}[Proof of Theorem~\ref{thm:1.1}]
\begin{EESegment}{EE-TGT-P124-001}
By \eqref{eq:10.5}, the field constructed above solves the equation exactly on
\NSi{NS-F-P124-001}{[0,1)}. Its smoothness and fixed compact support give
\NSi{NS-F-P124-002}{u\in C([0,T];H^3)} for every
\NSi{NS-F-P124-003}{T<1}. Lemma~\ref{lem:10.3} gives the force the required smoothness
and support, and Lemma~\ref{lem:10.4} gives the bounded energy.
\end{EESegment}

\begin{EESegment}{EE-TGT-P124-002}
Fix the \NSi{NS-F-P124-004}{X_{\mathrm{in}}} from Theorem~\ref{thm:3.1}, and take the
Cartesian points
\NSFormula{NS-F-P124-005}{display}{10.20}
\begin{equation}
  x_\tau=(\sqrt2X_{\mathrm{in}}\tau,0,0),\qquad t=1-\tau.
\tag{10.20}\label{eq:10.20}
\end{equation}
Along this path, \NSi{NS-F-P124-006}{z=0} and \NSi{NS-F-P124-007}{q=\tau}. For every
sufficiently small \NSi{NS-F-P124-008}{\tau>0}, the spatial and time cutoffs in
\eqref{eq:10.4} both equal one. Equation~\eqref{eq:3.6} then gives
\NSFormula{NS-F-P124-009}{display}{10.21}
\begin{equation}
  u_\theta(x_\tau,1-\tau)
  =\tau^{-A}\bigl(e_0+O(\tau^{2h})\bigr)\longrightarrow+\infty.
\tag{10.21}\label{eq:10.21}
\end{equation}
The annular corrections vanish in this fixed inner similarity region. In the expansion, the
cutoff sum of the terms with \NSi{NS-F-P124-010}{n\geq1}, in powers of
\NSi{NS-F-P124-011}{q^{2h}}, has exactly the relative-error bound shown here, with the
summation cutoffs still included. The points \NSi{NS-F-P124-012}{x_\tau} where the
velocity diverges approach the origin. Since
\NSi{NS-F-P124-013}{H^3(\R^3)\hookrightarrow L^\infty(\R^3)}, there can be no
classical \NSi{NS-F-P124-014}{H^3} continuation through time one. In the classical
\NSi{NS-F-P124-015}{H^3} class, Sobolev approximation justifies the same difference
energy calculation, and \NSi{NS-F-P124-016}{H^3\hookrightarrow W^{1,\infty}} supplies
the bound needed for Gronwall's inequality. This proves uniqueness on every shorter
interval. The maximal classical existence interval is therefore
\NSi{NS-F-P124-017}{[0,1)}.
\end{EESegment}

\begin{EESegment}{EE-TGT-P124-003}
Now suppose there were a global smooth solution \NSi{NS-F-P124-018}{v,P} with the same
data and uniformly bounded kinetic energy. For every \NSi{NS-F-P124-019}{T<1},
Lemma~\ref{lem:10.5} applies on the closed interval
\NSi{NS-F-P124-020}{[0,T]}, so
\NSi{NS-F-P124-021}{v=u} throughout \NSi{NS-F-P124-022}{[0,1)}. Equation
\eqref{eq:10.21} would then contradict the boundedness of the smooth field
\NSi{NS-F-P124-023}{v} on a compact neighborhood of
\NSi{NS-F-P124-024}{(0,1)}. The force cannot be zero, because if it were,
\eqref{eq:10.13} would give \NSi{NS-F-P124-025}{u=0}.
\end{EESegment}

\begin{EESegment}{EE-TGT-P124-004}
Finally, fix \NSi{NS-F-P124-026}{\nu>0} and define
\NSFormula{NS-F-P124-027}{display}{10.22}
\begin{equation}
\begin{aligned}
  u_\nu(x,t)&=\sqrt\nu\,u(x/\sqrt\nu,t),\\
  p_\nu(x,t)&=\nu p(x/\sqrt\nu,t),
  &f_\nu(x,t)&=\sqrt\nu\,f(x/\sqrt\nu,t).
\end{aligned}
\tag{10.22}\label{eq:10.22}
\end{equation}
Put \NSi{NS-F-P124-028}{y=x/\sqrt\nu}. The chain rule gives
\NSFormula{NS-F-P124-029}{display}{none}
\[
\begin{aligned}
  \partial_tu_\nu+(u_\nu\cdot\nabla_x)u_\nu-\nu\Delta_xu_\nu+\nabla_xp_\nu
  &=\sqrt\nu\,[\partial_tu+(u\cdot\nabla_y)u-\Delta_yu+\nabla_yp](y,t)
    =f_\nu(x,t),\\
  \nabla_x\cdot u_\nu(x,t)&=(\nabla_y\cdot u)(y,t)=0.
\end{aligned}
\]
The initial data are still zero, and the common spatial support becomes
\NSi{NS-F-P124-030}{\mathcal K_\nu=\sqrt\nu\mathcal K}. Time has not been rescaled,
and
\NSFormula{NS-F-P124-031}{display}{none}
\[
  \partial_x^\alpha\partial_t^m f_\nu(x,t)
  =\nu^{(1-|\alpha|)/2}
   (\partial_y^\alpha\partial_t^m f)(x/\sqrt\nu,t).
\]
So the force is still smooth and compactly supported in
\NSi{NS-F-P124-032}{\mathcal K_\nu\times[0,2]}, and it satisfies every required decay
bound for this fixed \NSi{NS-F-P124-033}{\nu>0}. The energy and dissipation scale as
\NSFormula{NS-F-P124-034}{display}{10.23}
\begin{equation}
  \|u_\nu(t)\|_2^2=\nu^{5/2}\|u(t)\|_2^2,\qquad
  \nu\int_0^1\|\nabla u_\nu\|_2^2\,dt
  =\nu^{5/2}\int_0^1\|\nabla u\|_2^2\,dt.
\tag{10.23}\label{eq:10.23}
\end{equation}
The growth occurs at \NSi{NS-F-P124-035}{\sqrt\nu x_\tau}, still at time one. If a
smooth bounded-energy solution \NSi{NS-F-P124-036}{v_\nu,P_\nu} existed for the
rescaled data, the inverse rescaling would give the viscosity-one solution
\NSFormula{NS-F-P124-037}{display}{none}
\[
  v(y,t)=\nu^{-1/2}v_\nu(\sqrt\nu y,t),\qquad
  P(y,t)=\nu^{-1}P_\nu(\sqrt\nu y,t),\qquad
  \|v(t)\|_2^2=\nu^{-5/2}\|v_\nu(t)\|_2^2.
\]
That viscosity-one solution has already been ruled out. This proves the theorem for every
fixed positive viscosity.
\end{EESegment}
\end{proof}

\NSPage{125}%
\begin{EESegment}{EE-TGT-P125-001}
\subsection{The periodic equations.}\label{sec:10.5}
The velocity, pressure, and force all have support in one fixed compact set
(Proposition~\ref{prop:10.1} and Lemma~\ref{lem:10.3}), so the construction can also be
made periodic. First, rescale the supports until they lie inside a fundamental cube. Then
add all their integer translates. The translated supports do not overlap, so the momentum
equation, including its nonlinear term, stays intact. The pressure becomes periodic too, as
the erratum to the problem statement requires \cite{ref-13}.

The periodic domain and the interior of its fundamental cube are
\NSFormula{NS-F-P125-001}{display}{none}
\[
  \T^3=\R^3/\Z^3,\qquad Q_0=(-1/2,1/2)^3.
\]
The next corollary proves the periodic breakdown alternative (D) in \cite{ref-13}.
\end{EESegment}

\begin{corollary}\label{cor:10.6}
\begin{EESegment}{EE-TGT-P125-002}
For every \NSi{NS-F-P125-002}{\nu>0}, there is a smooth force on
\NSi{NS-F-P125-003}{\T^3\times[0,\infty)} with compact support in time such that the
solution \NSi{NS-F-P125-004}{(U,P_{\mathrm{per}})} of the periodic Navier--Stokes
equations, with viscosity \NSi{NS-F-P125-005}{\nu} and zero initial velocity, is smooth
on \NSi{NS-F-P125-006}{[0,1)} and satisfies
\NSFormula{NS-F-P125-007}{display}{none}
\[
  \limsup_{t\uparrow1}\|U(t)\|_{L^\infty(\T^3)}=\infty.
\]
Before time one, its velocity and pressure are supported in a fixed compact subset of the
interior of \NSi{NS-F-P125-008}{Q_0}. No global smooth periodic solution exists for the
same initial data and force.
\end{EESegment}
\end{corollary}

\begin{proof}
\begin{EESegment}{EE-TGT-P125-003}
Fix the viscosity \NSi{NS-F-P125-009}{\nu}, and call the whole-space fields and force
already constructed \NSi{NS-F-P125-010}{u,p,f}. Their common spatial support lies in
\NSi{NS-F-P125-011}{\mathcal K_\nu=\sqrt\nu\mathcal K}, just as in the viscosity
rescaling \eqref{eq:10.22}. Choose \NSi{NS-F-P125-012}{\lambda>1} large enough that
\NSi{NS-F-P125-013}{\lambda^{-1}\mathcal K_\nu\Subset Q_0}, and set
\NSi{NS-F-P125-014}{t_0=1-\lambda^{-2}}. For \NSi{NS-F-P125-015}{t\geq t_0}, define
\NSFormula{NS-F-P125-016}{display}{none}
\[
\begin{aligned}
  \widetilde u(x,t)&=\lambda u(\lambda x,\lambda^2(t-t_0)),\\
  \widetilde p(x,t)&=\lambda^2p(\lambda x,\lambda^2(t-t_0)),\\
  \widetilde f(x,t)&=\lambda^3f(\lambda x,\lambda^2(t-t_0)).
\end{aligned}
\]
Use the first two formulas only for \NSi{NS-F-P125-017}{t<1}, and use the force formula
for every \NSi{NS-F-P125-018}{t\geq t_0}. Extend all three fields by zero when
\NSi{NS-F-P125-019}{t<t_0}. They are smooth across
\NSi{NS-F-P125-020}{t=t_0} because the original fields and force vanish on an interval
after their initial time. Every term in the momentum equation is
\NSi{NS-F-P125-021}{\lambda^3} times the corresponding original term, so the viscosity
is still \NSi{NS-F-P125-022}{\nu}. The original time one is reached exactly when
\NSi{NS-F-P125-023}{\lambda^2(t-t_0)=1}, which is
\NSi{NS-F-P125-024}{t=1}.
\end{EESegment}

\begin{EESegment}{EE-TGT-P125-004}
Define the periodic fields by
\NSFormula{NS-F-P125-025}{display}{none}
\[
\begin{aligned}
  U(x,t)&=\sum_{k\in\Z^3}\widetilde u(x+k,t),&&0\leq t<1,\\
  P_{\mathrm{per}}(x,t)&=\sum_{k\in\Z^3}\widetilde p(x+k,t),&&0\leq t<1,\\
  F_{\mathrm{per}}(x,t)&=\sum_{k\in\Z^3}\widetilde f(x+k,t),&&t\geq0.
\end{aligned}
\]
Different translates have disjoint supports, with a positive gap between them. These sums
are locally finite, so derivatives pass through them, and at most one translate is nonzero
at any point. In particular,
\NSFormula{NS-F-P125-026}{display}{none}
\[
  (U\cdot\nabla)U(x,t)
  =\sum_{k\in\Z^3}(\widetilde u(x+k,t)\cdot\nabla)\widetilde u(x+k,t).
\]
The periodic velocity, pressure, and force therefore solve the equation exactly:
\NSFormula{NS-F-P125-027}{display}{none}
\[
  \partial_tU+(U\cdot\nabla)U-\nu\Delta U+\nabla P_{\mathrm{per}}
  =F_{\mathrm{per}},\qquad
  \operatorname{div}U=0,\qquad U(\cdot,0)=0.
\]
\end{EESegment}

\NSPage{126}%
\begin{EESegment}{EE-TGT-P126-001}
The force is smooth, and its time support lies in
\NSi{NS-F-P126-001}{[t_0,1+\lambda^{-2}]}. On the torus, this gives every decay
estimate
\NSi{NS-F-P126-002}{\|\partial_x^\alpha\partial_t^mF_{\mathrm{per}}(t)\|_\infty
\leq C_{\alpha,m,N}(1+t)^{-N}}. For the path \NSi{NS-F-P126-003}{x_\tau} from
\eqref{eq:10.20}, the rescaled growth path is
\NSFormula{NS-F-P126-004}{display}{none}
\[
  \widetilde x_\tau=\lambda^{-1}\sqrt\nu x_\tau,\qquad
  \widetilde t_\tau=1-\lambda^{-2}\tau.
\]
This path stays inside \NSi{NS-F-P126-005}{Q_0} near
\NSi{NS-F-P126-006}{(0,1)}, where only the zeroth translate is nonzero. Equations
\eqref{eq:10.21} and \eqref{eq:10.22} give
\NSFormula{NS-F-P126-007}{display}{none}
\[
  U_\theta(\widetilde x_\tau,\widetilde t_\tau)
  =\lambda\sqrt\nu\,\tau^{-A}\bigl(e_0+O(\tau^{2h})\bigr)\longrightarrow+\infty.
\]
\end{EESegment}

\begin{EESegment}{EE-TGT-P126-002}
Finally, take any two smooth periodic solutions with the same force and initial data on an
interval \NSi{NS-F-P126-008}{[0,T]} with \NSi{NS-F-P126-009}{T<1}. Their difference
\NSi{NS-F-P126-010}{w} satisfies
\NSFormula{NS-F-P126-011}{display}{none}
\[
  \frac12\frac{d}{dt}\|w\|_{L^2(\T^3)}^2
  +\nu\|\nabla w\|_{L^2(\T^3)}^2
  \leq\|\nabla U\|_\infty\|w\|_{L^2(\T^3)}^2.
\]
Periodic integration makes both the transport term and the pressure term disappear.
Gronwall's inequality shows that the two solutions agree through
\NSi{NS-F-P126-012}{T}. Any global smooth solution would therefore equal
\NSi{NS-F-P126-013}{U} before time one, which is impossible because its velocity blows
up at \NSi{NS-F-P126-014}{t=1}.
\end{EESegment}
\end{proof}

\begin{EESegment}{EE-TGT-P126-003}
\appendix
\section{Matching radial moments and constructing the heat exterior}\label{sec:A}
This appendix constructs the outer part of the leading profile
\NSi{NS-F-P126-015}{(U,E,\Pi)} from Theorem~\ref{thm:4.6}. Its pressure integral gives
the datum \NSi{NS-F-P126-016}{\Pi_0} for the regular axis problem
(Lemma~\ref{lem:A.5}). Its moment conditions, including \eqref{eq:A.8}, make the
radially integrated leading residual stress \NSi{NS-F-P126-017}{\mathcal T_0} vanish in the
exterior after the heat replacement (Lemma~\ref{lem:A.8}). The construction has to meet
both requirements and still keep the strict stress cone inequalities on the chosen radial
intervals (Propositions~\ref{prop:A.4}, \ref{prop:A.7} and \ref{prop:A.10}).

It starts with finite-dimensional moment corrections (Lemmas~\ref{lem:A.1} and
\ref{lem:A.2}). The same corrections will later attach the axis profile
(Proposition~\ref{prop:B.8}) and change the shear (Proposition~\ref{prop:C.2}). Next, the
outer profile is built from a given inner reference profile (Section~\ref{sec:A.2}). At the
end, its terminal power law \NSi{NS-F-P126-018}{E_{\mathrm{pow}}} is replaced by an
exact heat flow. Compensating moment corrections keep both the pressure datum and the
exterior stress conditions unchanged (Lemmas~\ref{lem:A.6} and \ref{lem:A.8} and
Proposition~\ref{prop:A.7}).
\end{EESegment}

\begin{EESegment}{EE-TGT-P126-004}
\subsection{Adjusting finitely many moments.}\label{sec:A.1}
Changing a profile on a compact radial interval changes its cumulative moments in
\eqref{eq:4.15}. Add fixed smooth bumps with adjustable coefficients to correct those
changes. The next two results show why this works. Distinct power weights produce
independent changes in the moments (Lemma~\ref{lem:A.1}), and the resulting quadratic
moment equations can be solved when the discrepancies are small
(Lemma~\ref{lem:A.2}). Together with Lemma~\ref{lem:4.4}, these corrections let each
radial change preserve the data already prescribed farther out.
\end{EESegment}

\begin{lemma}\label{lem:A.1}
\begin{EESegment}{EE-TGT-P126-005}
Let \NSi{NS-F-P126-019}{\alpha_1,\ldots,\alpha_m} be distinct real numbers. For every
\NSi{NS-F-P126-020}{j}, let \NSi{NS-F-P126-021}{\beta_j} be a smooth function that is
nonnegative, not identically zero, and supported in the interior of a compact interval
\NSi{NS-F-P126-022}{I_j\subset(0,\infty)}, where
\NSi{NS-F-P126-023}{I_1<\cdots<I_m}. Then the matrix
\NSFormula{NS-F-P126-024}{display}{none}
\[
  B_{ij}=\int_0^\infty x^{\alpha_i}\beta_j(x)\,dx
\]
is invertible. The same statement holds in a coordinate obtained by taking a logarithm,
where the weights are exponential. If the data vary smoothly over a compact set of parameters and the stated
separations remain in place, then the inverse and every fixed parameter derivative of the
inverse stay bounded.
\end{EESegment}
\end{lemma}
